\section{Численные методы вычисления определителей и обратных матриц}
\label{sec:determinants-and-inverse-matrices}

Рассмотрим квадратную матрицу $A$, размерности $n \times n$. Поскольку при релизации многих численных методов возникает необходимость вычисления определителя матрицы или обратной матрицы, важно рассмотреть основные методы и их преимущества и недостатки.

\subsection{Вычисление определителя методом разложения по строкам/столбцам}

Для любой квадратной матрицы $A=(a_{ij})$:

\[
\det A=\sum_{j=1}^{n}a_{ij}C_{ij}
\]

— разложение по $i$-й строке, где

\[
C_{ij}=(-1)^{i+j}M_{ij},
\]

— алгебраическое дополнение,

а $M_{ij}$ — минор элемента $a_{ij}$. Минором $M_{ij}$ элемента $a_{ij}$ квадратной матрицы $A=(a_{ij})$ порядка $n$ называется определитель матрицы порядка $(n-1)$, полученной из матрицы $A$ вычёркиванием $i$-й строки и $j$-го столбца.

Можно также раскладывать по столбцу:

\[
\det A=\sum_{i=1}^{n}a_{ij}C_{ij}.
\]

Вычисление определителя матрицы $A$ разложением по строке или столбцу требует $O(n!)$ операций, поскольку на первом шаге задача сводится к $n$ аналогичным, но размерности $n-1$, и так далее. Очевидно, что вычисление $\det(A)$, например, для $n = 20$, потребует использования современного суперкомпьютера.

\subsection{Вычисление определителя методом Гаусса}

Используются преобразования:

\begin{enumerate}
	\item $R_i \leftrightarrow R_j$ — поменять строки местами; 
	\quad определитель меняет знак: $\det A \to -\det A$;
	
	\item $R_i \to kR_i$ — умножить строку на $k \neq 0$; 
	\quad определитель умножается на $k$: $\det A \to k\det A$;
	
	\item $R_i \to R_i + kR_j$ — прибавить к строке другую строку; 
	\quad определитель не изменяется: $\det A \to \det A$.
\end{enumerate}

\textbf{Условие применимости.} Все главные угловые миноры матрицы $A$ должны быть ненулевыми.\\

Приводим матрицу к треугольному виду:

\[
A \to
\begin{pmatrix}
	a_{11} & * & \cdots & * \\
	0 & a_{22} & \cdots & * \\
	\vdots & \vdots & \ddots & \vdots \\
	0 & 0 & \cdots & a_{nn}
\end{pmatrix}.
\]

После этого:

\[
\det A = a_{11}a_{22}\cdots a_{nn},
\]

с учётом изменений определителя при преобразованиях.

Сложность метода Гаусса (приведение к треугольному виду) для нахождения определителя составляет $O(n^3)$ арифметических операций.


\subsection{Вычисление определителя по $LU$-разложению}

\textbf{Теорема.} Пусть $A \in \mathbb{R}^{n \times n}$. Если все угловые миноры матрицы $A$ отличны от 0, то существует единственная нижняя унитреугольная матрица $L$ и единственная верхняя треугольная матрица $U$:

\[
L = 
\begin{pmatrix}
	1 & 0 & \cdots & 0 \\
	l_{21} & 1 & \cdots & 0 \\
	\vdots & \vdots & \ddots & \vdots \\
	l_{n1} & l_{n2} & \cdots & 1
\end{pmatrix}, \quad 
U = 
\begin{pmatrix}
	u_{11} & u_{12} & \cdots & u_{1n} \\
	0 & u_{22} & \cdots & u_{2n} \\
	\vdots & \vdots & \ddots & \vdots \\
	0 & 0 & \cdots & u_{nn}
\end{pmatrix},
\]

такие что $A = LU$.

Поскольку определители треугольных матриц равны произведению диагональных элементов,  
\[
\det A = \left( \prod_{i=1}^n l_{ii} \right) \left( \prod_{i=1}^n u_{ii} \right).
\]

Если $L$ имеет единицы на диагонали (так называемое $LU$-разложение без выбора главного элемента), то просто  
\[
\det A = \prod_{i=1}^n u_{ii}.
\]

Сложность $LU$-разложения для вычисления определителя составляет $O(n^3)$ арифметических операций.\\

\textbf{Замечание} Можно использовать и другие методы разложения, например, $QR$-разложение с учетом того, что $det(Q)=1$, а также другие подходящие разложения с учётом свойств матрицы $A$. 

\subsection{Вычисление обратной матрицы через алгебраические дополнения}

Матрица $A^{-1}$ называется обратной к квадратной матрице $A$, если
\[
A A^{-1} = A^{-1} A = E,
\]
где $E$ — единичная матрица.

Обратная матрица существует тогда и только тогда, когда $\det A \neq 0$ (матрица невырождена).\\

Если $\det A \neq 0$, то обратная матрица вычисляется по формуле
\[
A^{-1} = \frac{1}{\det A} \cdot \operatorname{adj}(A),
\]
где $\operatorname{adj}(A)$ — \textbf{присоединённая} (союзная) матрица, которая определяется как транспонированная матрица алгебраических дополнений.\\

Сложность вычисления обратной матрицы через присоединённую матрицу составляет $O(n \cdot n!)$.

\subsection{Вычисление обратной матрицы через решение систем (через $LU$ разложение)}

Обратная матрица $A^{-1}$ находится из уравнения
\[
A A^{-1} = E,
\]
где $E$ — единичная матрица.

Если $x^{(j)}$ — $j$-й столбец $A^{-1}$, а $e^{(j)}$ — $j$-й столбец $E$, то
\[
A x^{(j)} = e^{(j)}, \quad j = 1, 2, \dots, n.
\]

Таким образом, задача сводится к решению $n$ систем с одной матрицей $A$ и разными правыми частями.

\begin{enumerate}
	\item \textbf{LU-разложение:}
	\[
	A = LU,
	\]
	где $L$ — нижняя треугольная, $U$ — верхняя треугольная матрицы.
	
	\item \textbf{Для каждого} $j = 1, 2, \dots, n$:
	\[
	\begin{cases}
		L y^{(j)} = e^{(j)} & \text{(прямая подстановка)}, \\
		U x^{(j)} = y^{(j)} & \text{(обратная подстановка)}.
	\end{cases}
	\]
	
	\item \textbf{Формирование обратной матрицы:}
	\[
	A^{-1} = \begin{pmatrix}
		x^{(1)} & x^{(2)} & \cdots & x^{(n)}
	\end{pmatrix}.
	\]
\end{enumerate}


Сложность алгоритма:
\[
\underbrace{O(n^3)}_{\text{LU-разложение}} + \underbrace{n \cdot O(n^2)}_{\text{решение систем}} = O(n^3).
\]

\subsection{Метод Гаусса--Жордана}

Метод Гаусса--Жордана — это расширение метода Гаусса, которое позволяет найти обратную матрицу или решить СЛАУ за один проход, приводя матрицу $A$ не к треугольному, а к единичному виду $E$.

Для нахождения обратной матрицы составляется расширенная матрица:
\[
(A \mid E),
\]
и с помощью элементарных преобразований строк она приводится к виду
\[
(E \mid A^{-1}).
\]

Вычислительная сложность алгорима: $O(n^3)$.

\subsection{Обусловленность задачи решения СЛАУ}

Рассмотрим систему линейных уравнений
\[
Ax = b,
\]
где $A$ — невырожденная матрица.

Обусловленность характеризует чувствительность решения к малым изменениям входных данных. Число обусловленности матрицы $A$: $\operatorname{cond}(A) = \|A\| \cdot \|A^{-1}\|$. Для невырожденной матрицы $\operatorname{cond}(A) \geq 1$; если $\operatorname{cond}(A) \approx 1$ - задача хорошо обусловлена, если $\operatorname{cond}(A) >> 1$ - задача плохо обусловлена.\\

При наличии погрешностей и в $A$, и в $b$:
\[
\frac{\|\delta x\|}{\|x\|} \leq \frac{\operatorname{cond}(A)}{1 - \operatorname{cond}(A) \frac{\|\delta A\|}{\|A\|}} \left( \frac{\|\delta A\|}{\|A\|} + \frac{\|\delta b\|}{\|b\|} \right).
\]

При условии:
\[
\operatorname{cond}(A) \frac{\|\delta A\|}{\|A\|} < 1.
\]

\subsection{Численная устойчивость и выбор главного элемента}

Обусловленность относится к самой задаче, а устойчивость — к численному алгоритму. Плохо обусловленную задачу нельзя сделать хорошо обусловленной выбором алгоритма, но можно выбрать устойчивый алгоритм, который не будет вносить существенную дополнительную ошибку. \\

В качестве примера рассмотрим метод Жордана--Гаусса. Он может быть неустойчивым, если ведущий элемент близок к нулю. Для повышения устойчивости используют стратегии выбора главного элемента:
\textbf{Виды выбора главного элемента:}
\begin{enumerate}
	\item \textbf{Частичный (по столбцу):} 
	\[
	|a_{kj}| = \max_{i \geq j} |a_{ij}| \neq 0.
	\]
	На каждом шаге выбирается максимальный по модулю элемент в текущем столбце (среди строк $i \geq j$).
	
	\item \textbf{Полный (по всей подматрице):}
	\[
	|a_{kl}| = \max_{i \geq j, \, l \geq j} |a_{il}| \neq 0.
	\]
	Выбирается максимальный по модулю элемент во всей оставшейся подматрице.
\end{enumerate}

Это делается, чтобы избежать накопления ошибок округления ,избавиться от деленя на ноль или очень маленькое число.

\subsection{Что важно сказать на экзамене}

Нужно помнить основные методы вычисления определителя/обратной матрицы, понимать, преимущества/недостатки каждого метода, понимать, что такое обусловленность и устойчивость.

\newpage